\SourceLeaf{307}
\Needspace{8\baselineskip}\begingroup\small
\begin{center}
REMARKS.
\end{center}

The note \emph{Zur Theorie der Kegelschnitte} relates to the investigations mentioned in the letter of 12 May 1843. It is printed from the posthumous papers in the following division, Verwendung complexer Grössen für die Geometrie, on pp. 341 to 344 of this volume.

\textsc{Gauss}'s solution to the geometrical problem posed by \textsc{Möbius} was published by \textsc{Schumacher} in Astronomische Nachrichten, no. 42, 1823 November. It is reprinted in vol. IV, pp. 406---407.

Note [II] is on a separate leaf. It probably dates from the same time as the letter excerpts printed under [I]. \textsc{Gauss} applied the barycentric calculus in particular to the theory of the pentagon, as indicated under [3]. The developments in question, which are connected with the Pentagramma mirificum, can be found under the title \emph{Weitere Fragmente über das Pentagramma mirificum}, no. 11, pp. 109 to 111 of this volume (see also the accompanying remarks by \textsc{Fricke}, p. 117). See also the note \emph{Beziehung der Raumverhältnisse auf ein gegebenes Tetraeder}, vol. II, pp. 307 to 308. It probably dates from as early as July 1831.

\begin{flushright}
\textsc{Stäckel}.
\end{flushright}
\par\endgroup
